\section{Copying text}
%\label{sec:5:1:recopie}

In the program {\idt GreekCross\_1} below we find the same
sequence of commands four times, to draw the four
iden\-ti\-cal arms of the cross and turn. Each arm breaks down into two equal sides
of length 3 and an end of length 1 (the unit
being the length of a side of a square of the grid).

\vspace{\baselineskip}
\vspace*{3cm}%
\includevspacePostScript[-\baselineskip]{CroixGrecque}{0.95}{0.95}
\vspace*{-3cm}%
\label{prog:CroixGrecque_1}
\input python/GreekCross_1.tex

This situation, where the same sequence of commands of the \robot\ is found
several times in a row in a program, is very frequent because
geometric drawings are often made up of identical pieces. It is then
very tedious to type the same text several times. This can be avoided
by using the facilities offered by the computer for copying a text
already recorded in its memory and pasting it elsewhere in the
program. How to do this is explained for
PCs (under Windows or Linux) in figure \ref{fig:PC-copier-coller}
(page \pageref{fig:PC-copier-coller})
and for the Macintosh in figure \ref{fig:mac-copier-coller}
(page \pageref{fig:mac-copier-coller}).
\input COPIER-COLLER-PC.tex
\input COPIER-COLLER-MAC.tex

Making many copies by hand takes a long time. Moreover, one can easily imagine that
in a program several pages long, the reader has trouble finding the
identical parts of text. This phenomenon is illustrated in the
program {\idt Number888\_1} below. The difficulty would
obviously be even greater if the identical parts of text were
on different pages.

\vspace{\baselineskip}
\vspace*{2cm}%
\includevspacePostScript[-\baselineskip]{Nombre_888}{0.95}{0.95}
\vspace*{-2cm}%
\input python/Number888_1.tex

\section{Examples of procedures}
%\label{sec:5:2:ex-procedures}

Rather than copying texts by hand ourselves, we can
ask the \python\ interpreter to do it for us by using
{\em procedures}\index{Procedure}. Let us begin by rewriting the program
{\idt GreekCross\_1} (page \pageref{prog:CroixGrecque_1}) us\-ing procedures to bring
out the decomposition of the Greek cross into
four iden\-ti\-cal arms, each of them breaking down into two
equal sides of length 3 and an end of length 1.

\newpage
\vspace*{3cm}%
\includePostScript{CroixGrecque}{0.95}{0.95}
\vspace*{-3cm}%
\input python/GreekCross_2.tex

\noindent Using a procedure involves:
\begin{list}{$\scriptscriptstyle\bullet$}{
}
\item the {\em procedure declaration}\index{Procedure declaration}, which defines a sequence
of instructions (called the {\em body of the procedure})\index{Procedure body} and gives it a
name (called the {\em name of the procedure})\index{Identifier};
\item the {\em procedure call}\index{Procedure call}, which indicates that the sequence
of instructions mak\-ing up the body of the procedure must be copied
in place of the name of the procedure.
\end{list}

\noindent For the reader of the program, using procedures is
preferable to copying texts. Each time he or she sees the
name of a procedure, the reader understands that it stands for a copy of a text
that can be looked up in the procedure declaration. On meeting
several calls of the same procedure, the reader need make no effort
to understand immediately that it is the same text that is
copied. The program is therefore clearer. Isn't that the case for the
program {\idt Number888\_2} below, which is easier
to understand than the program {\idt Number888\_1}?

\vspace{\baselineskip}
\vspace*{2cm}%
\includevspacePostScript[-\baselineskip]{Nombre_888}{0.95}{0.95}
\vspace*{-2cm}%
\input python/Number888_2.tex

\section{Declaring procedures}
%\label{sec:5:3:decl-procedures}

The declarations of the procedures\index{Procedure declaration} of a program are placed between
the line `{\mcl from}\ {\idt LibRobot}\ {\mcl import}\ {\idt *}' and the instructions of the program.
The structure of a \python\ program is therefore the following:

\begin{center}
\begin{tabular}{@{}|p{12.5cm}|@{}} \hline
{\em \python\ program:}\\ \hline \hline
\begin{minipage}[t]{12.5cm}
{\medskip
\mbox{\blc\ \ \ {\mcl from}\hspace{0.3em}\blanc\hspace{0.3em}{\idt LibRobot}\hspace{0.3em}\blanc\hspace{0.3em}{\mcl import}\hspace{0.3em}\blanc\hspace{0.3em}{\idt *}}\par
\mbox{\blc\ \ \ \cmt \#\ \small\rm Comment explaining what the program does}\par
\mbox{\blc\ \ \ \small\rm  \ldots\ Declarations of the procedures of the program \ldots}\par
\mbox{\blc\ \ \ \small\rm  \ldots\ Instructions of the program%
\footnote{\ Including the calls of the procedures of the program.} (not indented) \ldots}\par
\medskip}
\end{minipage}\\ \hline
\end{tabular}
\end{center}


\noindent In \python, a procedure declaration has the following form
(we shall use more elaborate forms later on):

\begin{center}
\begin{tabular}{@{}|p{12.5cm}|@{}} \hline
{\em Procedure declaration:}\\ \hline \hline
\begin{minipage}[t]{12.5cm}
{\medskip
\mbox{\blc\ \ \ \cmt \#\ \small\rm Specification: what the procedure does}\par
\mbox{\blc\ \ \ {\mcl def}\hspace{0.3em}\blanc\hspace{0.3em}\idt procedureName():}\par
\mbox{\blc\ \ \ \blk\blk\blk\blk\small\rm \ldots\ Body of the procedure%
\footnote{\ Including the calls of previously declared procedures.}
 ({\bf necessarily indented}) \ldots}\par
\medskip}
\end{minipage}\\ \hline
\end{tabular}
\end{center}

The English word {\mcl def} is short for {\em define}.
The {\idt  procedureName} is chosen by the programmer
as he or she likes, so as to suggest what
this procedure is for. The parentheses {\idt ()} that follow it indicate
that it is a procedure, which will have to be called by writing {\idt procedureName()},
exactly like the commands of the \robot\ {\idt fd()}, {\idt qtr()}, etc.
Indeed, the commands of the \robot\ are procedures declared in the
library {\idt LibRobot}. The {\mcl def} line must end with
a colon `{\idt :}', which announces the body of the procedure.

\index{Indentation}The instructions of the body of the procedure are written on the following
lines, all shifted to the right by the same number of blanks,
usually 4. They are said to be {\em indented}.
As in the program, one or more instructions can be written per line, separated
by semicolons. It is the indentation, and the indentation alone, that tells \python\
which instructions are part of the body of the procedure: the body ends at
the first line that is no longer indented. Blanks at the beginning of a line are therefore
not simply a matter of layout: they are part of the syntax of \python.
If you forget to indent the first line of the body, \python\ reports for example
`{\tt IndentationError: expected an indented block after function definition on line 3}'.
If on the contrary you indent an instruction of the program by mistake, it reports
`{\tt IndentationError: unexpected indent}'. And if a line of the
body does not have the same number of blanks as the previous ones, it reports
`{\tt IndentationError: unindent does not match any outer indentation level}'.
An empty line after the body
of each procedure, as in the program {\idt Number888\_2} above,
is not compulsory but makes the end of the procedure clearly visible.

{\bf Warning:} the body of a procedure {\bf must} be indented, and
{\bf all} its lines by the same number of blanks. A line of the body that you forget
to indent is not always reported as an error: it is simply no longer
part of the procedure. In the example below, {\idt qtr()} is not indented:
\python\ executes it only once, at the beginning of the program, and not at each call of
{\idt corner()}. Instead of a square, the \robot\ draws a straight line eight steps long, without
any error message:

{\parindent \parindentprog
\medskip
\noindent\begin{minipage}{\linewidth}\parindent \parindentprog
\vspace{\neglskip}\mbox{\blc {\mcl def}{\idt \ corner():}}\par
\vspace{\neglskip}\mbox{\blc \blk\blk{\idt fd();\ fd()}}\par
\vspace{\neglskip}\mbox{\blc {\idt qtr()}\ \ \ \ \ \ \ \ \ \ \ {\cmt \#\ Error:\ not\ indented,\ outside\ the\ procedure}}\par
\vspace{0.3\baselineskip}
\vspace{\neglskip}\mbox{\blc {\idt corner();\ corner();\ corner();\ corner()}}\par

\end{minipage}
\medskip }

\noindent For the \robot\ to turn at each corner, {\idt qtr()} must be indented
like {\idt fd();\ fd()}, with the same four blanks.

\medskip
\noindent {\it Remark:} in \python, procedures are special cases of
what are called {\em functions}\index{Function}. We shall come back to them
later.


\section{Identifiers}
%\label{sec:7:2:Ident}
The name of a procedure is called an
{\em identifier}\index{Identifier} in \python. Here is
a precise definition, to which you should
refer every time we speak of iden\-ti\-fi\-ers:

\begin{center}
\begin{tabular}{@{}|p{12.5cm}|@{}} \hline
{\em Identifier:}\\ \hline \hline
\begin{minipage}[t]{12.5cm}
{\medskip
An {\em identifier} is a sequence of unaccented lowercase
or uppercase letters, digits or
underscores `\_', which does not begin with a digit and which is different from the
English words reserved by \python\ such as {\mcl from}, {\mcl import}, {\mcl def}, {\mcl for},
{\mcl in}, {\mcl if}, {\mcl while}, {\mcl return}, {\mcl True}, etc.
\medskip}
\end{minipage}\\ \hline
\end{tabular}
\end{center}

For example {\idt side}, {\idt greekCross}, {\idt greek\_cross} or {\idt side2} are
identifiers whereas {\idt 2side}, {\idt left-side} or {\idt for} are not.
If you write `{\mcl def}\ {\idt 2side():}', \python\ answers
`{\tt SyntaxError: invalid decimal literal}' because it thinks it is
reading a number; for `{\mcl def}\ {\idt left-side():}' or `{\mcl def}\ {\idt for():}', it
simply answers `{\tt SyntaxError: invalid syntax}'.
Uppercase and lowercase letters are distinguished: {\idt Side} and {\idt side} are two
different identifiers. If you declare `{\mcl def}\ {\idt Side():}' and then
call {\idt side()}, \python\ answers
`{\tt NameError: name \textquotesingle{}side\textquotesingle{} is not defined. Did you mean: \textquotesingle{}Side\textquotesingle{}?}'.
\python\ allows an identifier to begin with an uppercase letter, but the custom is to
begin the names of procedures with a lowercase letter. To make
an identifier made up of several words readable, we
put an uppercase letter at the beginning of each word except the first, as in
{\idt saintAndrewsCross} or {\idt innerEdgeOfLens}\footnote{\ Many \python\
programmers prefer to separate the words with underscores, as in {\idt saint\_andrews\_cross}.}.

\section{Calling procedures}
%\label{sec:5:4:app-procedures}

\index{Procedure call}To call a procedure, you just write its name followed
by {\idt ()}. The call {\idt procedureName()}
is then replaced by the body of the procedure\index{Procedure body}, that is, by the list
of indented instructions that follows the line `{\mcl def}\ {\idt procedureName():}' in its declaration.

Note that if you declare a procedure without ever calling it, it
is of absolutely no use.

A procedure can be called in the program or in another procedure.
For example in the program {\idt GreekCross\_2}, we call
the procedure {\idt cross}, which calls the procedure {\idt arm}, which itself
calls the procedures {\idt side} and {\idt end}.

The rule that a procedure must
be
de\-clared before being called must absolutely be followed. Suppose for example that we have
declared a pro\-ce\-dure under the name {\idt procName}. Suppose also
that, because of a typing mistake, we have called this procedure
under the name {\idt progName}. When execution reaches this call, the interpreter
stops, reporting the error
`{\tt NameError: name \textquotesingle{}progName\textquotesingle{} is not defined}'.
Not knowing that
we wrote {\idt progName} by mistake instead of {\idt procName}, the interpreter
assumes that the procedure {\idt procName} has been declared but has never been used.
This is useless but allowed. It also notes that the procedure
{\idt progName} is called but has not been declared. This is forbidden.
It notices, however, that the two names look alike and suggests
`{\tt Did you mean: \textquotesingle{}procName\textquotesingle{}?}'.%
\footnote{\ We often speak of programs in an
anthropomorphic way, that is, by attributing to them characteristics
proper to human beings. In fact we are thinking of the programmers who
wrote them. For example, the computer does not make mistakes; it is its
manufacturer who made a mistake, or the programmer who made a logical
error in his or her program.}
We get the same `{\tt NameError}' error if the program calls a procedure
whose declaration is placed further on: \python\ executes the declarations in
the order in which they are written, and at the time of the call it does not yet know the procedure.
Like all errors discovered during execution, these errors appear
only once the \robot\ has made the part of the drawing that precedes the faulty call.

\section{Robot commands defined by procedures}
%\label{sec:5:5:robot-cmd-procs}

Some commands of the \robot\ are equivalent to a sequence
of elementary commands and can therefore be expressed with procedures.
For example, to turn a quarter turn, it is enough to turn an eighth of a turn twice
in a row. We therefore have:
\input python/Pivotements1.tex
\noindent Here are new commands of the \robot\ that can be expressed by means of
procedures using commands we already know:
\input python/Pivotements2.tex
\needspace{6\baselineskip}\noindent \rule{137mm}{0.1mm}\par\nopagebreak %-------------------------------------------------------
\begin{exercices}{
Write \python\ programs with pro\-ce\-dures to
draw the crosses,
%de la figure \ref{fig:Croix},
the astroid star, the magnifying glass or the notched double square
%de la figure \ref{fig:EtLoCa}
below.
\noindent $\Box$}\end{exercices}
\begin{figure}[htbp]
 \vspace{-\baselineskip}
 \includevspacePostScript[1.7cm]{Croix}{0.95}{0.95}
 \caption{Cross potent, double cross, papal cross, cross botonny and Saint Andrew's cross}
 %\label{fig:Croix}
\end{figure}
\begin{figure}[htbp]
 \includevspacePostScript[5.1cm]{EtoileLoupeCarre}{0.95}{0.95}
 \caption{Astroid star, magnifying glass and notched double square}
 %\label{fig:EtLoCa}
\end{figure}


\begin{corrige}{
We give one example of a program ({\idt AstroidStar}), the
others being solved in the folder {\idt book/\allowbreak 05\_procedures}.
\input python/AstroidStar.tex
\noindent $\Box$}\end{corrige}

\vfill

%%%%%%%%%%%%%%%%%%%%%%%%%%%%%%%%%%%%%%%%%%%%%%%%%%%%%%%%%%%%%%%%%%%%%%%%%%%
\noindent \framebox[\textwidth]{
\begin{minipage}{\myparboxwidth}
\setlength{\parindent}{2em}
\small \sf
\vspace{\abovedisplayskip}
\section*{Addenda}
The notion of procedure\index{Procedure} introduced by Von Neumann and Turing
was used in 1949
under the name of {\em routine} on the EDSAC ({\em Electronic Delay
Storage Automatic Computer}), generally considered
the first stored-program computer to have
worked. This machine was built from 1947 at the University
of Cambridge in England by the team of Wilkes\footnote{\ Maurice
Wilkes, one of the inventors of radar.} and ran its first programs
in May 1949 \cite{levy1}.

Procedures illustrate the concept of {\em abstraction}\index{Abstraction}: the
properties of an object are isolated so as to
use them or study them separately. More precisely, a
procedure gives a name to an algorithm for constructing part
of a drawing. In the body of the procedure,
we are interested in how to construct the drawing. In the calls
we are interested only in the finished drawing, while wishing to ignore completely how
the drawing is obtained. We thus separate the object from the way of
constructing it, the what from the how.

\vspace{\abovedisplayskip}
\end{minipage}
}

%%%%%%%%%%%%%%%%%%%%%%%%%%%%%%%%%%%%%%%%%%%%%%%%%%%%%%%%%%%%%%%%%%%%%%%%%%%%
%%%%%%%%%%%%%%%%%%%%%%%%%%%%%%%%%%%%%%%%%%%%%%%%%%%%%%%%%%%%%%%%%%%%%%%%%%%
\noindent \framebox[\textwidth]{
\begin{minipage}{\myparboxwidth}
\setlength{\parindent}{2em}
\small \sf
\vspace{\abovedisplayskip}
We find it useful to describe the drawing made by a procedure
in a comment
placed at the beginning of the procedure, called the
 {\em specification}\index{Specification} of the procedure. Note that for a procedure
to be usable without ever looking at its body, the
specification of the procedure must describe which drawing it makes but
also what the positions and orientations
of the \robot\ are before and after the drawing is made. This brings out
a notion of {\em invariant}\index{Invariant}, to which we shall return in the following
chapters on iteration.

Procedures are the basis of the methods
for structuring very large programs: to be
understandable, programs of several thousand lines
must be split into procedures no longer than
a few pages. Procedures are therefore an elegant means of
dividing a program up logically. However, the concept of procedure
is not sufficient to structure programs of several million
lines. One then uses the notion of {\em module} (also called {\em unit})
and of {\em library} to group together sets of procedures.

Procedures grouped together in a
{\em library} (in \python\ also called a {\em module}) can be
supplied to programmers to help them in their tasks. To use
a procedure, the programmer must know its name and its
specification, without needing
to know how it works internally. The line `{\mcl from}\ {\idt LibRobot}\ {\mcl import}\ {\idt *}' at the beginning
of the program is there to tell the interpreter which
library is used in this program. The commands of the \robot\ are
procedures grouped together in libraries.

Analysing geometric drawings to bring out a structure that
can be described algorithmically by means of procedures is very
good training for the mind. In this sense learning to program
can play for children a formative role similar to the one
once played by Euclidean geometry.

\vspace{\abovedisplayskip}
\end{minipage}
}
%%%%%%%%%%%%%%%%%%%%%%%%%%%%%%%%%%%%%%%%%%%%%%%%%%%%%%%%%%%%%%%%%%%%%%%%%%%%
\newpage
\thispagestyle{empty}
